\chapter{Implementación}
\label{ch:implementacion}
% El desarrollo software; no volcar código entero. Fuente: src/ + README.
% Frontera con metodología: aquí solo el CÓMO; el QUÉ/POR QUÉ está en \ref{ch:metodologia}.
% No repetir decisiones de diseño; describir su realización.
% Plantilla de capítulo: abrir con un párrafo de hoja de ruta que enumere estas secciones con
% \ref, y cerrar con \section{Conclusiones del capítulo}, al modo de \ref{sec:fund-conclusiones}.
% Ambas se escriben con el cuerpo, no antes: un título sin prosa no debe llegar al PDF.

\section{Arquitectura del pipeline}
\label{sec:pipeline}
% src/: config, data, models, metrics, metrics_runner, logger, train; logging a parquet.

\section{Cálculo de las métricas}
\label{sec:impl-metricas}
% Interfaz común compute(model,X,y,loss_fn) sobre el batch de medición; REGISTRY+BASELINE; last-layer en ResNet.

\section{Reproducibilidad}
\label{sec:reproducibilidad}
% Principio: seed.py, determinismo como requisito, normalización por dataset, caveat Tiny-ImageNet
% (~6e-4, libjpeg/Pillow). El detalle reproducible (entorno, versiones, hardware) va en el anexo
% \ref{sec:determinismo}; remitir, no repetir.

\section{Ejecución y reanudación}
\label{sec:ejecucion}
% run_matrix.py idempotente (hecho = summary.json); troceable por dataset/model/optimizer.
% OJO: una sola GPU, no hay cluster (ver \ref{sec:motivacion}). El troceado sirve para reanudar
% por tramos, no para repartir entre nodos.

\section{Pilot de calibración}
\label{sec:pilot}
% run_pilot.py: 1 run/celda, presupuesto doblado, reports_pilot/. Aquí solo la realización
% software (aislamiento de salidas, reanudación); el protocolo de calibración y los valores
% finales están en \ref{sec:calibracion}.
